\documentclass[10pt,a4paper]{amsart}
\usepackage{amsmath,amssymb,amsthm,graphicx}
% Retained minimal preamble for the Stacks Project chapter varieties.tex.
% Verbatim excerpt of the chapter preamble: Xy-pic setup (native preamble lines 14-18), the native theorem
% declarations with the mathematical macros (native preamble lines 154-217), and the project own no-op
% definitions of the reference/slogan/history display environments (native wrapper.tex lines 1-3) that the
% counted statement uses for its slogan. Copied mechanically; nothing added, removed or redefined.
% For commutative diagrams we use Xy-pic
\usepackage[all]{xy}

% We use 2cell for 2-commutative diagrams.
\xyoption{2cell}
% Theorem environments.
%
\theoremstyle{plain}
\newtheorem{theorem}[subsection]{Theorem}
\newtheorem{proposition}[subsection]{Proposition}
\newtheorem{lemma}[subsection]{Lemma}

\theoremstyle{definition}
\newtheorem{definition}[subsection]{Definition}
\newtheorem{example}[subsection]{Example}
\newtheorem{exercise}[subsection]{Exercise}
\newtheorem{situation}[subsection]{Situation}

\theoremstyle{remark}
\newtheorem{remark}[subsection]{Remark}
\newtheorem{remarks}[subsection]{Remarks}

\numberwithin{equation}{subsection}

% Macros
%
\def\lim{\mathop{\mathrm{lim}}\nolimits}
\def\colim{\mathop{\mathrm{colim}}\nolimits}
\def\Spec{\mathop{\mathrm{Spec}}}
\def\Hom{\mathop{\mathrm{Hom}}\nolimits}
\def\Ext{\mathop{\mathrm{Ext}}\nolimits}
\def\SheafHom{\mathop{\mathcal{H}\!\mathit{om}}\nolimits}
\def\SheafExt{\mathop{\mathcal{E}\!\mathit{xt}}\nolimits}
\def\Sch{\mathit{Sch}}
\def\Mor{\mathop{\mathrm{Mor}}\nolimits}
\def\Ob{\mathop{\mathrm{Ob}}\nolimits}
\def\Sh{\mathop{\mathit{Sh}}\nolimits}
\def\NL{\mathop{N\!L}\nolimits}
\def\CH{\mathop{\mathrm{CH}}\nolimits}
\def\proetale{{pro\text{-}\acute{e}tale}}
\def\etale{{\acute{e}tale}}
\def\QCoh{\mathit{QCoh}}
\def\Ker{\mathop{\mathrm{Ker}}}
\def\Im{\mathop{\mathrm{Im}}}
\def\Coker{\mathop{\mathrm{Coker}}}
\def\Coim{\mathop{\mathrm{Coim}}}

% Boxtimes
%
\DeclareMathSymbol{\boxtimes}{\mathbin}{AMSa}{"02}

%
% Macros for moduli stacks/spaces
%
\def\QCohstack{\mathcal{QC}\!\mathit{oh}}
\def\Cohstack{\mathcal{C}\!\mathit{oh}}
\def\Spacesstack{\mathcal{S}\!\mathit{paces}}
\def\Quotfunctor{\mathrm{Quot}}
\def\Hilbfunctor{\mathrm{Hilb}}
\def\Curvesstack{\mathcal{C}\!\mathit{urves}}
\def\Polarizedstack{\mathcal{P}\!\mathit{olarized}}
\def\Complexesstack{\mathcal{C}\!\mathit{omplexes}}
% \Pic is the operator that assigns to X its picard group, usage \Pic(X)
% \Picardstack_{X/B} denotes the Picard stack of X over B
% \Picardfunctor_{X/B} denotes the Picard functor of X over B
\def\Pic{\mathop{\mathrm{Pic}}\nolimits}
\def\Picardstack{\mathcal{P}\!\mathit{ic}}
\def\Picardfunctor{\mathrm{Pic}}
\def\Deformationcategory{\mathcal{D}\!\mathit{ef}}
\newenvironment{reference}{}{}
\newenvironment{slogan}{}{}
\newenvironment{history}{}{}
\usepackage{cleveref}
\crefname{originallemma}{original lemma}{original lemmas}
\Crefname{originallemma}{Original lemma}{Original lemmas}
\makeatletter
\newlabel{lemma-inverse-image-irreducible}{{04KW}{source}{}{sourcepointer-1}{}}
\newlabel{lemma-inverse-image-irreducible@cref}{{[originallemma][0][]04KW}{[1][0][]source}{}{}{}}
\makeatother
\crefname{originallemma}{original lemma}{original lemmas}
\Crefname{originallemma}{Original lemma}{Original lemmas}
\makeatletter
\newlabel{lemma-separably-closed-field-irreducible-components}{{038H}{source}{}{sourcepointer-2}{}}
\newlabel{lemma-separably-closed-field-irreducible-components@cref}{{[originallemma][0][]038H}{[1][0][]source}{}{}{}}
\makeatother
\crefname{originallemma}{original lemma}{original lemmas}
\Crefname{originallemma}{Original lemma}{Original lemmas}
\makeatletter
\newlabel{morphisms-lemma-generalizations-lift-flat}{{03HV}{source}{}{sourcepointer-3}{}}
\newlabel{morphisms-lemma-generalizations-lift-flat@cref}{{[originallemma][0][]03HV}{[1][0][]source}{}{}{}}
\makeatother
\crefname{originallemma}{original lemma}{original lemmas}
\Crefname{originallemma}{Original lemma}{Original lemmas}
\makeatletter
\newlabel{morphisms-lemma-integral-universally-closed}{{01WM}{source}{}{sourcepointer-4}{}}
\newlabel{morphisms-lemma-integral-universally-closed@cref}{{[originallemma][0][]01WM}{[1][0][]source}{}{}{}}
\makeatother
\title{Irreducible components of a base change map onto irreducible components, and the map is surjective}
\author{The Stacks Project Authors}
\date{stacks-project:varieties.tex:local-snapshot-20261011; GNU FDL 1.2 or later}
\begin{document}
\maketitle
\noindent\textit{Editorial note.} All delivered mathematical text is the author's own, copied from the pinned Stacks Project chapter file varieties.tex; no statement or proof step was rewritten, reordered or removed.\par\medskip
\noindent\textit{Editorial note.} The proof is complete as written and ends with the surjectivity argument; no step is referred to the reader.\par\medskip
\noindent\textit{Editorial note.} The counted spans carry 4 navigation stubs and no bibliography citation.\par\medskip
\noindent\textit{Editorial note.} The preamble is the chapter's retained minimal preamble (sources/stacks-varieties/retained-preamble.tex lines 1-77), a verbatim excerpt of the native preamble.\par\medskip
\noindent\textit{Editorial note.} The label lemma-image-irreducible is defined only in varieties.tex, so no cross-chapter restatement had to be weighed for this card.\par\medskip
\noindent\textit{Allowed standard prerequisites (editorial).}\begin{itemize}
\item Base change along separable algebraic closures preserves irreducible components. Same chapter, Lemma lemma-separably-closed-field-irreducible-components, invoked at proof lines 1645-1648. Supplies the reduction to K equal to the separable closure of k.
\item Generalizations lift along flat morphisms. Morphisms, Lemma morphisms-lemma-generalizations-lift-flat, invoked at proof lines 1655-1656. Supplies that generic points map to generic points.
\item Integral morphisms are universally closed. Morphisms, Lemma morphisms-lemma-integral-universally-closed, invoked at proof lines 1659-1660. Supplies that the image of an irreducible component is closed.
\item Inverse image of an irreducible component is a nonempty union of components. Same chapter, Lemma lemma-inverse-image-irreducible, invoked at proof lines 1665-1667. Supplies the surjectivity statement.
\end{itemize}\medskip
\ifdefined\hypertarget\hypertarget{sourcepointer-1}{}\fi
\noindent\textit{Source reference.} Original lemma 04KW: sources/stacks-varieties/pinned-original-source.tex, source lines 1596--1596. Pointer outside the counted span, invoked once for context; the load-bearing argument of the counted statement is the author's proof reproduced in full in the proof span.\par\medskip
\ifdefined\hypertarget\hypertarget{sourcepointer-2}{}\fi
\noindent\textit{Source reference.} Original lemma 038H: sources/stacks-varieties/pinned-original-source.tex, source lines 1524--1524. Pointer outside the counted span, invoked once for context; the load-bearing argument of the counted statement is the author's proof reproduced in full in the proof span.\par\medskip
\ifdefined\hypertarget\hypertarget{sourcepointer-3}{}\fi
\noindent\textit{Source reference.} Original lemma 03HV: sources/stacks-chow/tags.txt, source lines 4581--4581. Pointer outside the counted span, invoked once for context; the load-bearing argument of the counted statement is the author's proof reproduced in full in the proof span.\par\medskip
\ifdefined\hypertarget\hypertarget{sourcepointer-4}{}\fi
\noindent\textit{Source reference.} Original lemma 01WM: sources/stacks-chow/tags.txt, source lines 2487--2487. Pointer outside the counted span, invoked once for context; the load-bearing argument of the counted statement is the author's proof reproduced in full in the proof span.\par\medskip
\begin{lemma}
\label{lemma-image-irreducible}
Let $K/k$ be an extension of fields.
Let $X$ be a scheme over $k$.
For every irreducible component $\overline{T} \subset X_K$
the image of $\overline{T}$ in $X$ is an irreducible component in $X$.
This defines a canonical map
$$
\text{IrredComp}(X_K)
\longrightarrow
\text{IrredComp}(X)
$$
which is surjective.
\end{lemma}

\begin{proof}
Consider the diagram
$$
\xymatrix{
X_K \ar[d] & X_{\overline{K}} \ar[d] \ar[l] \\
X & X_{\overline{k}} \ar[l]
}
$$
where $\overline{K}$ is the separable algebraic closure of $K$, and
where $\overline{k}$ is the separable algebraic closure of $k$. By
Lemma \ref{lemma-separably-closed-field-irreducible-components}
the morphism $X_{\overline{K}} \to X_{\overline{k}}$ induces
a bijection between irreducible components. Hence it suffices
to show the lemma for the morphisms
$X_{\overline{k}} \to X$ and $X_{\overline{K}} \to X_K$.
In other words we may assume that $K = \overline{k}$.

\medskip\noindent
The morphism $p : X_{\overline{k}} \to X$ is integral, flat and surjective.
Flatness implies that generalizations lift along $p$, see
Morphisms, Lemma \ref{morphisms-lemma-generalizations-lift-flat}.
Hence generic points of irreducible components of $X_{\overline{k}}$
map to generic points of irreducible components of $X$.
Integrality implies that $p$ is universally closed, see
Morphisms, Lemma \ref{morphisms-lemma-integral-universally-closed}.
Hence we conclude that the image $p(\overline{T})$ of an irreducible component
is a closed irreducible subset which contains a generic point of an
irreducible component of $X$, hence $p(\overline{T})$
is an irreducible component of $X$. This proves the first assertion.
If $T \subset X$ is an irreducible component, then $p^{-1}(T) =T_K$
is a nonempty union of irreducible components, see
Lemma \ref{lemma-inverse-image-irreducible}.
Each of these necessarily maps onto $T$ by the first part.
Hence the map is surjective.
\end{proof}

\end{document}
